\begin{abstract}
On-device Retrieval-Augmented Generation (RAG) enables privacy-preserving, low-latency question answering on resource-constrained mobile hardware. However, executing autoregressive Small Language Models (SLMs) on mobile CPUs incurs severe prompt-prefill latency when repeatedly encoding retrieved context. While precomputing and restoring Key-Value (KV) cache prefixes can eliminate redundant prefill computation, existing mobile runtimes treat KV-cache files as opaque conversational histories. This introduces a critical systems failure mode: \emph{silent stale-cache hallucination}, where outdated or revoked retrieved evidence preserved in serialized KV caches causes the model to generate factually incorrect responses without detection. Furthermore, dynamic context-capacity constraints and continuation-token boundary effects frequently trigger out-of-memory (OOM) aborts or silent truncation on memory-limited handsets.

In this paper, we present \textbf{HACER-RAG} (\textbf{H}ierarchical \textbf{A}daptive \textbf{C}ache, \textbf{E}vidence-version-aware \textbf{R}epair for on-device RAG), a mobile inference system designed for reliable and efficient local RAG on commodity Android smartphones. HACER-RAG pairs a canonical evidence layout with cryptographic chunk-level hashing to bind every KV cache segment to an explicit evidence version. When retrieved context changes, HACER-RAG executes a selective repair policy that preserves valid prefix KV tensors while triggering online prefill only for updated downstream chunks, falling back safely to clean scratch prefill under cache corruption. In addition, HACER-RAG introduces a JNI-level context-budget safety guard that dynamically enforces headroom constraints and models continuation boundary-token alignment. 

We evaluate HACER-RAG on a physical Xiaomi Redmi 12 5G smartphone (Qualcomm Snapdragon 4 Gen 2 SoC, Android 14) using a 4-bit quantized Gemma 2 2B model executing via \texttt{llama.rn} and \texttt{llama.cpp}. In a rigorous 90-matched-pair empirical scaling study across nine configurations ($K \in \{1, 2, 4\}$ chunks, $n_{\text{ctx}} \in \{512, 1024, 2048\}$), HACER-RAG achieves mean prompt-prefill compute reductions of \textbf{19.3--21.2\%} for $K=1$ (84 cached tokens), \textbf{38.5--40.4\%} for $K=2$ (158 cached tokens), and \textbf{86.1--86.7\%} for $K=4$ (341 cached tokens), delivering up to \textbf{3.65$\times$} end-to-end wall-clock speedup with zero native truncation and 100\% deterministic output parity. A balanced Two-Way ANOVA demonstrates that prefill reductions are strictly governed by cached token count ($F = 1920.56, p < 10^{-68}, \eta_p^2 = 0.979$), with no statistically significant main effect from context capacity ($p = 0.981$) or interaction ($p = 0.729$). In repeated multi-chunk RAG trials (9 matched pairs), HACER-RAG achieves a mean prefill reduction of \textbf{44.13\%} and a wall speedup of \textbf{1.63$\times$} while preserving 100\% multi-hop factual accuracy, incurring an average flash cache-load overhead of only 256.0~ms compared to 16.8~seconds of CPU compute saved. Across five dynamic corpus failure scenarios (content updates, document revocation, chunk reordering, partial updates, and cache corruption), HACER-RAG completely eliminates stale-cache hallucinations and runtime crashes where naive session reuse fails.
\end{abstract}

\begin{IEEEkeywords}
On-Device Machine Learning, Retrieval-Augmented Generation, Key-Value Cache Reuse, Mobile Systems, Small Language Models, Cache Invalidation, Latency Optimization.
\end{IEEEkeywords}
